% Source: https://artofproblemsolving.com/community/c3h2251003_dice_and_bars
% Stable identity: aops-topic-2251003-q01
There are $6^3=216$ equally likely outcomes, and $\dfrac{1}{8}=\dfrac{27}{216}$, so we need the sums that occur in exactly $27$ ways.

\textbf{Sum $10$.} The number of solutions of $a+b+c=10$ with $1\le a,b,c\le6$: putting $a=a'+1$ etc., $a'+b'+c'=7$ with $0\le a',b',c'\le5$. Without the upper bound there are $\dbinom92=36$ solutions; those with some variable at least $6$ have that variable equal to $6$ or $7$ and the other two summing to $1$ or $0$: $3\cdot(2+1)=9$. So sum $10$ occurs in $36-9=27$ ways.

\textbf{Symmetry.} Replacing each die value $v$ by $7-v$ turns a sum $s$ into $21-s$, so sums $s$ and $21-s$ occur equally often; hence sum $11$ also occurs $27$ times.

\textbf{No other sums.} The numbers of ways for the sums $3,4,\dots,10$ are
\[1,3,6,10,15,21,25,27,\]
increasing, and by symmetry the sums $11,\dots,18$ give $27,25,\dots,1$. So exactly the sums $10$ and $11$ have probability $\dfrac{1}{8}$, and the answer is $10+11=\boxed{21}$.
